\documentclass[a4paper,12pt]{article}
\usepackage[utf8]{inputenc}
\usepackage[T1]{fontenc}
\usepackage[ngerman]{babel}
\usepackage{amsmath}
\usepackage{amssymb}
\usepackage{enumitem}
\usepackage{geometry}
\geometry{a4paper, margin=2.5cm}
\usepackage{parskip}

\title{Einsendeaufgaben -- Aufgabe 5.5\\[0.5em]
\large Modul 61112: Lineare Algebra}
\author{}
\date{}

\begin{document}
\maketitle

\section*{Aufgabe 5.5}

\begin{enumerate}[label=\alph*)]
  \item
  Nach dem Fundamentalsatz der Algebra zerfällt das charakteristische Polynom
  $\chi_\psi \in \mathbb{C}[X]$ in Linearfaktoren. $\psi$ besitzt also
  mindestens einen Eigenwert $\lambda$ mit dem Eigenraum
  $\operatorname{Eig}_\psi(\lambda)$, da $V \neq \{0\}$.

  Wir erkennen, dass für alle $v \in \operatorname{Eig}_\psi(\lambda)$ und für
  alle $\varphi \in F$ gilt:
  \begin{align*}
    \psi(\varphi(v)) &= (\psi \circ \varphi)(v) \\
                     &= (\varphi \circ \psi)(v)
                       && \text{(wegen Voraussetzung } \psi \circ \varphi = \varphi \circ \psi\text{)} \\
                     &= \varphi(\psi(v)) \\
                     &= \varphi(\lambda v)
                       && \text{(wegen } v \in \operatorname{Eig}_\psi(\lambda)\text{)} \\
                     &= \lambda \varphi(v)
                       && \text{(wegen Linearität von } \varphi\text{)}
  \end{align*}
  Demnach ist $\varphi(v) \in \operatorname{Eig}_\psi(\lambda)$.

  Daraus können wir schließen, dass für jedes $\varphi$ der Eigenraum
  $\operatorname{Eig}_\psi(\lambda)$ $\varphi$-invariant ist, da
  \[
    \varphi(v) \in \operatorname{Eig}_\psi(\lambda) \quad \text{für alle }
    v \in \operatorname{Eig}_\psi(\lambda).
  \]

  Da $\dim(\operatorname{Eig}_\psi(\lambda)) \geq 1$, folgt
  $\{0\} \subsetneq \operatorname{Eig}_\psi(\lambda)$.

  Außerdem ist noch $\operatorname{Eig}_\psi(\lambda) \subsetneq V$ zu zeigen.
  Damit $\operatorname{Eig}_\psi(\lambda) = V$ muss $\psi - \lambda \operatorname{id}_V$
  die Nullabbildung sein, denn nur dann gilt $\operatorname{Ker}(\psi - \lambda \operatorname{id}_V) = V$.
  $\psi - \lambda \operatorname{id}_V = 0_V$ ist äquivalent zu
  $\psi = \lambda \operatorname{id}_V$. $\psi = \lambda \operatorname{id}_V$ ist
  allerdings als Voraussetzung ausgeschlossen, weswegen nach Kontraposition
  $\operatorname{Eig}_\psi(\lambda) \subsetneq V$ folgt.

  Zusammengefasst gilt also
  $\{0\} \subsetneq \operatorname{Eig}_\psi(\lambda) \subsetneq V$ sowie
  $\operatorname{Eig}_\psi(\lambda)$ $\varphi$-invariant für alle
  $\varphi \in F$, was zu zeigen war.

  \item
  \textbf{$\subseteq$-Richtung: $Z \subseteq \{\lambda \operatorname{id}_V \mid \lambda \in \mathbb{C}\}$}

  Wir zeigen $Z \subseteq \{\lambda \operatorname{id}_V \mid \lambda \in \mathbb{C}\}$
  mit Hilfe der Kontraposition von a) mit $F := \operatorname{End}(V)$.

  Wir müssen also zeigen, dass für alle Unterräume $U$ von $V$ gilt:
  \begin{align*}
    &\text{Entweder } U = \{0\} \\
    &\text{oder } U = V \\
    &\text{oder es existiert ein } \varphi \in \operatorname{End}(V), \\
    &\text{sodass } U \text{ nicht } \varphi\text{-invariant ist.}
  \end{align*}

  Dann können wir die Kontraposition von a) anwenden, woraus folgt, dass für
  alle $\psi \in Z$
  \[
    \psi = \lambda \operatorname{id}_V \quad \text{für ein } \lambda \in \mathbb{C} \text{ gilt,}
  \]
  also $\psi \in \{\lambda \operatorname{id}_V \mid \lambda \in \mathbb{C}\}$.

  Sei $U$ ein beliebiger Unterraum von $V$. Wenn $U = \{0\}$ oder $U = V$ ist
  nichts zu tun. Wir können also voraussetzen $\{0\} \subsetneq U \subsetneq V$.
  Wir wählen eine beliebige Basis $B_U$ aus. Da $U \subsetneq V$, muss es
  mindestens ein $v$ geben, um $B_U$ zu ergänzen zu einer Basis $B$ von $V$.
  Wir definieren eine lineare Abbildung $\varphi \in \operatorname{End}(V)$
  durch die Abbildung der Basis $B = \{b_1, \ldots, b_n\}$:
  \[
    \varphi(b_i) = v \quad \text{für alle } 1 \leq i \leq n.
  \]
  Da $v \notin U$, ist $U$ nicht $\varphi$-invariant.

  Wir können also die Kontraposition anwenden und erhalten
  $Z \subseteq \{\lambda \operatorname{id}_V \mid \lambda \in \mathbb{C}\}$.

  \textbf{$\supseteq$-Richtung: $\{\lambda \operatorname{id}_V \mid \lambda \in \mathbb{C}\} \subseteq Z$}

  Sei $\lambda \in \mathbb{C}$ beliebig. Es folgt direkt
  \[
    \varphi \circ \lambda \operatorname{id}_V = \lambda \varphi
    = \lambda \operatorname{id}_V \circ \varphi
    \quad \text{für alle } \varphi \in \operatorname{End}(V).
  \]
  Also $\lambda \operatorname{id}_V \in Z$ und dadurch
  $\{\lambda \operatorname{id}_V \mid \lambda \in \mathbb{C}\} \subseteq Z$.
\end{enumerate}

\end{document}
